\DocumentMetadata{
  pdfversion=2.0,
  pdfstandard=ua-2,
  lang=en-US,
  tagging=on,
  tagging-setup={math/setup=mathml-SE,tabsorder=structure}
}
\documentclass[11pt,letterpaper]{article}
\usepackage[margin=0.68in,includefoot]{geometry}
\usepackage{newcomputermodern}
\usepackage{amsmath,amscd,mathtools}
\usepackage{tikz,adjustbox}
\providecommand{\square}{\mdlgwhtsquare}
\usepackage[hidelinks]{hyperref}
\hypersetup{
  pdftitle={Analysis PhD Exam, May 2017},
  pdfauthor={University of Florida Department of Mathematics},
  pdfsubject={Graduate mathematics examination},
  pdfdisplaydoctitle=true,
  bookmarksopen=true
}
\setcounter{secnumdepth}{0}
\setlength{\parindent}{0pt}
\setlength{\parskip}{0.28em}
\setlength{\emergencystretch}{3em}
\setlength{\leftmargini}{2.15em}
\setlength{\leftmarginii}{2.65em}
\setlength{\leftmarginiii}{3em}
\renewcommand{\labelenumi}{\textbf{\theenumi.}}
\pagestyle{empty}
\raggedbottom
\allowdisplaybreaks
\newenvironment{examproblems}[1]{%
  \begin{enumerate}%
  \setcounter{enumi}{#1}%
  \setlength{\topsep}{0.35em}%
  \setlength{\partopsep}{0pt}%
  \setlength{\itemsep}{0.48em}%
  \setlength{\parsep}{0.18em}%
}{\end{enumerate}}
\newenvironment{examsubparts}{%
  \begin{enumerate}%
  \setlength{\topsep}{0.18em}%
  \setlength{\partopsep}{0pt}%
  \setlength{\itemsep}{0.16em}%
  \setlength{\parsep}{0.1em}%
}{\end{enumerate}}
\newenvironment{examinstructions}{%
  \par\begingroup\itshape
}{%
  \par\endgroup\vspace{0.18em}
}
\makeatletter
\renewcommand\section{\@startsection{section}{1}{0pt}%
  {-1.25ex plus -.25ex minus -.1ex}%
  {0.55ex plus .1ex}%
  {\normalfont\large\bfseries}}
\renewcommand{\@maketitle}{%
  \newpage\null\vskip -1em%
  {\footnotesize\bfseries
    \href{https://www.ufl.edu/}{University of Florida}, \href{https://math.ufl.edu/}{Department of Mathematics}\par}%
  \vskip -0.5em%
  \tagstructbegin{tag=Title}%
    {\LARGE\bfseries\href{https://gma.math.ufl.edu/exams/analysis-phd/}{Analysis PhD Exam}, May 2017\par}%
  \tagstructend%
  \vskip 0.25em%
  \tagmcbegin{artifact}%
    \hrule height 1pt%
  \tagmcend%
  \vskip 1em%
}
\makeatother
\title{Analysis PhD Exam, May 2017}
\author{}
\date{}
\begin{document}
\maketitle
\thispagestyle{empty}
\begin{examinstructions}
Do six of eight. Problems (2) and (5) are short answer. For the remaining problems, please provide complete proofs with all details.
\end{examinstructions}
\begin{examproblems}{0}
\item
Prove that if \(X,Y\) are topological spaces and \(f:X\to Y\) is continuous, then~\(f\) is Borel measurable.
\item
Do three of four.
\begin{examsubparts}
\item[\textnormal{(a)}]
Give an example, if possible, of a closed set \(E\subset\mathbb R\) of positive Lebesgue measure that contains no interval.
\item[\textnormal{(b)}]
Give an example, if possible, of a sequence \(f_n:[0,1]\to\mathbb R\) of Lebesgue measurable functions that converges in measure, but not pointwise a.e.
\item[\textnormal{(c)}]
Give an example, if possible, of a premeasure on a Boolean algebra \(\mathcal A\) without a unique extension to a measure on the~\(\sigma\)-algebra generated by \(\mathcal A\).
\item[\textnormal{(d)}]
Determine the limit of the sequence \((a_n)_{n=1}^{\infty}\), defined by 
\[
a_n=\int_1^{\infty}\frac{\cos^2(\pi t)}{t^n}\,dt.
\]
\end{examsubparts}
\item
Suppose \((X,\mathcal M,\mu)\) is a~\(\sigma\)-finite measure space and \(\mathcal N\subset\mathcal M\) is a (sub)~\(\sigma\)-algebra. Let \(\nu=\mu|_{\mathcal N}\). Show, if \((X,\mathcal N,\nu)\) is a~\(\sigma\)-finite measure space and \(f\in L^1(\mu)\), then there exists a \(g\in L^1(\nu)\) such that, for all \(E\in\mathcal N\), 
\[
\int_E g\,d\nu=\int_E f\,d\mu.
\]
 Is the~\(\sigma\)-finiteness hypothesis on~\(\nu\) necessary?
\item
Given \(f\in L^1([0,1])\), let \(g(x)=\int_x^1\frac{f(t)}{t}\,dt\). Show \(g\in L^1([0,1])\) and 
\[
\int_0^1 g=\int_0^1 f.
\]
\item
Do three of four.
\begin{examsubparts}
\item[\textnormal{(a)}]
Is there a norm on the vector space 
\[
c_{00}=\{f:\mathbb N\to\mathbb C:\text{ there exists an }N\text{ such that }f(n)=0\text{ for }n\ge N\}
\]
 that makes \(c_{00}\) a Banach space?
\item[\textnormal{(b)}]
Explain what is meant by the Fourier transform of a function \(f\in L^2(\mathbb R)\).
\item[\textnormal{(c)}]
Is there bounded linear bijection between \(C([0,1])\) and \(L^\infty([0,1])\)?
\item[\textnormal{(d)}]
Give an example, if possible, of a sequence of unit vectors \((f_n)\) from \(L^2(\mathbb R)\) that converges weakly to some \(f\in L^2(\mathbb R)\), but not in norm.
\end{examsubparts}
\item
Show, if \(E\subset\mathbb R\) is a Lebesgue measurable set of positive measure, then \(E-E=\{x-y:x,y\in E\}\) contains an open interval.
\item
Suppose \(\mathcal X\) is a Banach space and \(\mathcal M\) and \(\mathcal N\) are closed subspaces. Show, if for each \(x\in\mathcal X\) there exist unique \(m\in\mathcal M\) and \(n\in\mathcal N\) such that 
\[
x=m+n,
\]
 then the mapping \(P:\mathcal X\to\mathcal M\) defined by \(Px=m\) is bounded.
\item
Suppose \((X,\mathcal M,\mu)\) is a~\(\sigma\)-finite measure space and \(1\le q<p<\infty\). Show, \(L^q\subset L^p\) if and only if there exists a \(\delta>0\) such that if \(E\in\mathcal M\) then either \(\mu(E)=0\) or \(\mu(E)\ge\delta\).
\end{examproblems}
\end{document}
